%% kapa: exam=1 round=26 year=2015 period=1 form=A subject=civil_law subject_ko=민법 seq=39 number=39
%% source: Q-Net 감정평가사 기출문제 (cst00309) artlSeq=5138001 file=extracted/1008/5138001/제26회_감정평가사 1차_기출/제26회_감정평가사 1차_기출/1교시_감평_A형.pdf
%% conversion: kordoc 4.24.51 -> tools/gen1.py (auto); review=proofread-sonnet; flags=none
\begin{question}{39}
유치권에 관한 설명으로 옳은 것은? (다툼이 있으면 판례에 따름)
\begin{choices}
  \item 채권자가 불법으로 점유를 취득한 경우에도 유치권이 성립한다.
  \item 채권자가 유치권을 행사하면 채권의 소멸시효는 중단된다.
  \item 건물임차인은 권리금반환청구권에 기하여 임차건물에 대하여 유치권을 주장할 수 없다.
  \item 유치권에는 우선변제적 효력이 없으므로, 유치권자는 채권의 변제를 받기 위하여 유치물을 경매할 수 없다.
  \item 공사대금채권에 기하여 유치권을 행사하는 자가 채무자의 승낙 없이 유치물의 보존에 필요한 범위 내에서 유치물인 주택에 거주하며 사용하였다면, 특별한 사정이 없는 한 차임에 상당한 이득을 소유자에게 반환할 의무가 없다.
\end{choices}
\end{question}
\answer{3}
